\begin{proposition}[the actual cusp-link inclusion]\label{push:src:C.3.16}
Orient the central orbit cycle $H=W/S^1_\delta\subset Q=X_c/S^1_\delta$
by the complex $q$-normal orientation, where $q=1/t$ near the cusp.
Use the marked positive period $\gamma$ and the actual transverse
$q$-meridian of \ref{push:src:C.3.13}. Then
\[
 H_1(Q\setminus H;\Z)=\Z,
 \qquad \operatorname{lk}(\gamma,H)=12.
\]
More precisely, write the cusp orbit link as
$L=S^1_\gamma\times\operatorname{Nil}^3(u,w,q)$, with the integral
basis $(\gamma,u,q)$ of $H_1(L;\Z)$ fixed in \ref{push:src:C.3.13}. The inclusion
homomorphism is
\[
 H_1(L;\Z)\longrightarrow H_1(Q\setminus H;\Z),
 \qquad (a,b,c)\longmapsto 12a+c,
\]
where the target generator is the positive meridian of $H$.
This is a statement about the actual cusp inclusion, not only the
weighted target three-sphere.
\end{proposition}
\Uses{the hypotheses, notation and definitions stated here; \cref{push:src:C.3.13}; \cref{push:src:C.3.14}; \cref{push:src:C.3.7}}

\begin{proof}
By \ref{push:src:C.3.14}, $Q$ is the standard smooth five-sphere. The finite
normalization pushout in \ref{push:src:C.3.7} makes $H$ a compact, locally
contractible space, and gives $H^3(H;\Z)=\Z$. Its oriented
normalization has degree one at a generic smooth point; identification
along the conductor, of dimension at most two, does not change the
top-dimensional fundamental cycle. Thus $H$ has the multiplicity-one
oriented fundamental cycle specified in the statement. Integral
Alexander duality gives
\[
 \widetilde H_1(Q\setminus H;\Z)
 \simeq \widetilde H^3(H;\Z)=\Z.
\]
The local positive meridian at any generic smooth point pairs to one
with that fundamental cycle and therefore represents the indicated
primitive generator.

Let $\overline G:Q\to S^3$ be the continuous descent of the smooth
$S^1_\delta$-invariant map constructed in \ref{push:src:C.3.13}, and write its two
complex coordinate functions as $G_1,G_2$. Define the globally
single-valued continuous complex function
\[
 h=G_2^3-G_1^4.
\]
Since the weighted target map is
$p(z_1,z_2)=z_2^3/(z_2^3-z_1^4)$ and
$p\circ\overline G=\overline f$, the zero set of $h$ is exactly $H$.
In the cusp solid-torus formula of \ref{push:src:C.3.13}, it is explicitly
\[
 h=q\,S(q)^{12}\exp(24\pi i\vartheta),\qquad
 \vartheta=g-\operatorname{Re}\kappa_0(q),
\]
where $S(q)>0$, $g$ is circle-valued, and $\kappa_0$ is a
single-valued holomorphic function on the cusp disc. The factor
multiplying $q$ is a continuous nowhere-zero function. Hence on a
small positive normal circle at a generic smooth point of $H$, the
map $h/|h|$ has degree $+1$. It follows that
\[
 [h/|h|]\in H^1(Q\setminus H;\Z)
\]
is exactly the integral meridian, or linking, class. This argument
requires only continuity of the descended map at fixed points; it
does not assert that an arbitrary invariant smooth map descends
smoothly in every orbit-space coordinate system.

On the outer cusp collar the extension $g$ equals the raw marked
character $\theta$. Traversing the positive $\gamma$ period changes
$\theta$ by one, while $q$ and the other torus coordinates stay fixed.
The displayed formula for $h$ consequently has argument degree
$12$ on that loop. A traversal of $u$ changes neither $q$ nor
$\theta$, and has degree zero. Finally, the actual positive
transverse meridian is the section $z_1=z_2=1$ of \ref{push:src:C.3.13}. On it
$\theta=0$ identically, so
\[
 \frac{h}{|h|}
 =\frac{q}{|q|}\exp\bigl(-24\pi i\operatorname{Re}\kappa_0(q)\bigr).
\]
The second factor extends over the whole disc and has degree zero.
The meridian therefore has degree one. These three evaluations give
the asserted inclusion homomorphism. The commutator direction $w$
is already zero in $H_1(L;\Z)$.
\end{proof}

\begin{remark}[Scope]
The numerical value $12$ is now an actual linking invariant in the
orbit five-sphere. The function $q$ alone cannot provide it: globally
$q$ is a meromorphic base coordinate with poles, so its argument
compares the zero and pole fibres. The globally defined function $h$
removes that ambiguity. The present result specifies one part of the
global attaching data for the cap of \ref{push:src:C.3.14}; it does not evaluate its
Pontryagin framing term or its characteristic integer $J_\psi$.
\end{remark}
\Uses{the hypotheses, notation and definitions stated here; \cref{push:src:C.3.14}}


